\section{Study design}\label{sec:method}

The model of Section~\ref{sec:model} is silent on the sign of the first-round window; this section fixes the protocol that measures it. The baseline is C0 and the intervention is FRO-2, both defined in Section~\ref{sec:intro}.

For each variant within each condition, we estimate the DFR as the number of failed decodings divided by the number attempted. Interval uncertainty comes from a $95\%$ bootstrap over keys rather than over error draws, so intervals reflect key-to-key variation. A zero-failure entry is never read as zero risk; the rule-of-three bound $3/N$ replaces it over $N$ attempts. Every comparison reduces to the risk ratio, C0's DFR divided by the variant's, and values above one favour the variant. Variants differ only in the threshold schedule; the iteration budget and every other setting are held fixed.

Interval and test conventions were fixed before the runs. Intervals are percentile cluster-bootstrap intervals over keys: 10{,}000 resamples per contrast, keys drawn with replacement while every error draw of a selected key is retained, with the bootstrap seed fixed at 20260905. The sign test counts the keys on which the variant fails less often than C0, drops ties, and returns an exact two-sided binomial probability. Holm correction covers the three formal conditions of the primary comparison; the temporal grid, the transfer checks and the decoder comparisons are reported as marginal intervals. All computations are reproducible from the archived per-key counts.

The study uses eight operating points, each with a role. The primary point is $P_1'=(4801,45,92)$; the prime marks the measurable variant of the screening point $P_1=(4801,45,84)$ at which candidate schedules were compared before formal measurement. The transfer points are $P_2=(3001,37,66)$ and $P_3=(4001,41,82)$, and $P_4=(3607,39,76)$ is held out for a rule check. The largest block in the series is $(6007,56,100)$. At the production block $(12323,71)$ the study measures a raised weight $(t^{\ast}=172)$ and the production weight 134; the CPA-128 preset is $(10163,71,140)$.

Three key/error sampling conditions are used throughout. RAND draws the key at random and the error uniformly at weight $t$. TYPE-I constrains one parity-check block of the key to contain a run of $12$ consecutive ones, with the error drawn as in RAND; it is an adversarial probe, not a population estimate. NCW-shallow uses a key pool filtered to admit a near-codeword pattern and builds the error with overlap $\kappa=12$, below the block row weight $d=45$. The shallow condition is a stress construction, not a proof that capture is impossible; a separate stress family that crosses $d$ appears in Section~\ref{sec:boundary}. Figure~\ref{fig:design} shows the pairing and the run streams.

The protocol ran in stages, from calibration to a single confirmation on fresh keys. We first calibrated the primary family where failures are measurable, so every compared difference sits above the zero-failure floor. Before prespecification we screened nine candidate schedules under a worst-group bound, which reads the weakest condition so a schedule cannot hide harm in one condition behind gains in another. The prespecified gate required the worst-group $95\%$ upper bound to improve by at least $0.5$ in $\log_{10}$ units (a factor of 3.2), and the primary condition had to stay non-inferior. The gate was evaluated once on a pure-fresh subset: 212 keys with 3000 draws each, 636{,}000 per variant. Those keys never entered screening, so the confirmation is independent of every selection decision; the subset remains part of the 278-key formal pool, so this independence is from the selection procedure, not a second statistically independent estimate. Four frozen master seeds govern the master, bootstrap, exploratory and confirmation streams, so every run is exactly reproducible.

Every number reported later is reproducible from the frozen exports. A second implementation, written independently within the same project, recorded 11 failures in 20{,}000 dynamic-distance decodes against 136 in 200{,}000 in the C reference ($z=-0.68$), and reproduced the direction of the primary contrast. That cross-check bounds coding error only, because both implementations share the design assumptions. An exact-recovery audit covered all 278 formal $P_1'$ keys and 834{,}000 decodings per variant. It found 795 exact-recovery failures for C0 and 167 for FRO-2, matching the zero-syndrome totals, with no miscorrection in either variant. The main experiment uses static-distance semantics while the C reference uses dynamic distances, so their rates are never compared. Exploratory entries and archived prototypes remain labelled and never carry a primary conclusion.

\dblfig{fig1_study_design.pdf}{Study design. (a)~Both schedules decode the same key, error and syndrome. (b)~FRO-2 adds two in round one only; from round two onward both use the C0 schedule. (c)~Outcomes are paired per key, and keys are the resampling unit.}{fig:design}
